\documentclass[11pt]{article}
\usepackage[utf8]{inputenc}
\usepackage{amsmath}
\usepackage{tikz}
\usetikzlibrary{shapes,arrows,positioning}
\usepackage{graphicx}
\begin{document}
NAME-AYUSH MISHRA\\
ENROLLMENT NUMBER -0801CS211029  
$$JEE ADVANCED PAPER$$\\ 
 $$PART-1 PHYSICS$$\\
  $$SECTION1$$\\
\begin{center}
\begin{tikzpicture}
\node[draw]{\textbullet The section contains 2 question each question for +3 marks};
\end{tikzpicture} \\
\end{center}       
Q.1 A heavy nucleus Q of half-life 20 minutes undergoes alpha-decay with probability of 60\% and beta-decay with probability of 40\%. Initially, the number of Q nuclei is 1000. The number of alpha-decays of Q in the first one hour is\\
(A)50\\    (B)75 \\      (C)350 \\     (D)525\\
Q.2 A football of radius R is kept on a hole of radius r (r< R) made on a plank kept horizontally. One end of the plank is now lifted so that it gets tilted making an angle $\theta$ from the horizontal as shown in the figure below. The maximum value of $\theta$ so that the football does not start rolling down the plank satisfies (figure is schematic and not drawn to scale)
\newpage
\begin{center}
\begin{figure}
\includegraphics[scale=0.5]{physicsimage1.png};
\end{figure}
\end{center}

(A)$\sin\theta=\frac{r}{R}$\\
(B)$\tan\theta=\frac{r}{R}$\\
(C)$\sin\theta=\frac{r}{2R}$\\
(D)$\cos\theta=\frac{r}{2R}$\\
                        $$SECTION2$$\\
\begin{center}
\begin{tikzpicture}
\node[draw]{\textbullet The section contains 2 question each question for +4 marks};
\end{tikzpicture} \\
\end{center}  
Q.3 A thin rod of mass M and length $\alpha$ is free to rotate in horizontal plane about a fixed vertical axis passing through point O. A thin circular disc of mass M and of radius $\alpha /4$ is pivoted on this rod with its center at a distance $\alpha /4$ from the free end so that it can rotate freely about its vertical axis, as shown in the figure. Assume that both the rod and the disc have uniform density and they remain horizontal during the motion. An outside stationary observer finds the rod rotating with an angular velocity  $\omega$ and the disc rotating about its vertical axis with angular velocity $4\omega$ The total angular momentum of the system about the point O is $\frac{Ma^2\omega}{48}$. The value of n is .
\newpage
\begin{center}
\begin{figure}
\includegraphics[scale=0.5]{physicsimage2.png};
\end{figure}
\end{center}
Q.4 Two large circular discs separated by a distance of 0.01m are connected to a battery via a switch as shown in the figure. Charged oil drops of density $900 kg m^3$ are released through a tiny hole at the center of the top disc. Once some oil drops achieve terminal velocity, the switch is closed to apply a voltage of 200V across the discs. As a result, an oil drop of radius $8×10^7$ m stops moving
vertically and floats between the discs. The number of electrons present in this oil drop is ?(neglect the buoyancy force, take acceleration due to gravity = 10 m $s^{-2}$ and the charge on the electron (e) = 1.6*$10^{-19}C$)
\newpage
\begin{center}
\begin{figure}
\includegraphics[scale=0.5]{physicsimage3.png};
\end{figure}
\end{center}

                      $$SECTION3$$\\                                      
Q.5 Two conducting cylinders of equal length but different radii are connected in series between two heat baths kept at temperatures $T_1=300K$and $T_2=100$ 𝐾, as shown in the figure.The radius of the bigger cylinder is twice that of the smaller one and the thermal conductivities of the materials of the smaller and the larger cylinders are $k_1$ and $k_2$ respectively. If the temperature at the junction of the two cylinders in the steady state is 200K,then $\frac{k_1}{k_2}$=

\newpage
\begin{center}
\begin{figure}
\includegraphics[scale=0.5]{physics4.png}
\end{figure}
\end{center}



                        $$CHEMISTRY$$\\
                        $$SECTION1$$\\
\begin{center}
\begin{tikzpicture}
\node[draw]{\textbullet The section contains 2 question each question for +3 marks};
\end{tikzpicture} \\
\end{center} 
Q.6 A colorless aqueous solution contains nitrates of two metals, X and Y. When it was added to an aqueous solution of NaCl, a white precipitate was formed. This precipitate was found to be partly soluble in hot water to give a residue P and a solution Q. The residue P was soluble in aq. $NH_3$ and also in excess sodium thiosulfate. The hot solution Q gave a yellow precipitate with KI. The metals X and Y, respectively, are
(A) Ag and Pb \\
(B) Ag and Cd\\
(C) Cd and Pb\\
(D) Cd and Zn\\                     
Q.7 Which among the following statement(s) is (are) true for the extraction of aluminium from bauxite?\\
(A) Hydrated $Al_2 O_3$ precipitates, when $CO_2$ is bubbled through a solution of sodium\\
(B) Addition of $Na_3 AlF_6$ lowers the melting point of alumina.\\
(C) $CO_2$ is evolved at the anode during electrolysis.\\
(D) The cathode is a steel vessel with a lining of carbon.\\ 
\newpage
                          $$SECTION2$$\\
\begin{center}
\begin{tikzpicture}
\node[draw]{\textbullet The section contains 2 question each question for +4 marks};
\end{tikzpicture} \\
\end{center}
Q.8 The $1^{rd},2^{nd} and 3^{rd}$ ionization enthalpies,$I_1,I_2 and I_3$of four atoms with atomic numbers n, n + 1, n + 2, and n + 3 , where n < 10, are tabulated below. What is the value of n? 
\begin{center}
\begin{tabular}{|c|c|c|c|}
\hline
atomic number & $I_1$ & $I_2$ & $I_3$\\\hline
n & 1681 & 3374 & 6050\\\hline
n+1 & 2081 & 3952 &6122\\\hline
n+2 & 496 & 4562 & 6910\\\hline
n+3 & 738 & 1451 & 7733\\\hline
\end{tabular}
\end{center}
Q.9 Consider the Bohr's model of a one-electron atom where the electron moves around the nucleus. In the following, List-I contains some quantities for the $n_{th}$ orbit of the atom and list-II contains options showing how they depend on n\\
\begin{tabular}{c c}
list-I & list-II\\
(I)Radius of the $n^{th}$ orbit & (P)$\alpha n^{-2}$\\
(II)Angular momentum of the electron in the $n^th$ orbit & (Q)$\alpha n^{-1}$\\
(III)Kinetic energy of the electron in the $n^th$ orbit & (R) $\alpha n^0$ \\
(IV)Potential energy of the electron in the $n^{th}$ orbit & (S)$\alpha n^1$\\
 & (T) $\alpha n^2$\\
 & (U) $\alpha n^{\frac{1}{2}}$\\
\end{tabular}
Which of the following options has the correct combination considering List-I and List-II?\\
(A) (I),(P) \\
(B) (I),(T)\\
(C) (II),(Q)\\
(D) (II),(R)\\
\newpage
                    $$SECTION3$$\\
\begin{center}
\begin{tikzpicture}
\node[draw]{\textbullet The section contains 2 question each question for +3 marks};
\end{tikzpicture} \\
\end{center}

Q.10 The cubic unit cell structure of a compound containing cation M and anion X is shown below. When compared to the anion, the cation has smaller ionic radius. Choose the correct statement(s).
\newpage
\begin{center}
\begin{figure}
\includegraphics[width=0.5\textwidth]{foruseinlatex.png}
\end{figure}
\end{center}

(A)	The empirical formula of the compound is MX. \\
(B)	The cation M and anion X have different coordination geometries.\\ 
(C)	The ratio of M-X bond length to the cubic unit cell edge length is 0.866.\\ 
(D)	The ratio of the ionic radii of cation M to anion X is 0.414.


                     $$MATHEMATICS$$
                      $$SECTION1$$\\
\begin{center}
\begin{tikzpicture}
\node[draw]{\textbullet The section contains 2 question each question for +3 marks};
\end{tikzpicture} \\
\end{center}
Q.11 The probability that a missile hits a target successfully is 0.75. In order to destroy the target completely, at least three successful hits are required. Then the minimum number of missiles that have to be fired so that the probability of completely destroying the target is NOT less than 0.95, is 
\begin{equation}
(A)x^2+y^2-3x+y=0
\end{equation}
\begin{equation}
(B)x^2+y^2+x+3y=0
\end{equation}
\begin{equation}
(C)x^2+y^2+2x-1=0
\end{equation}
\begin{equation}
(D)x^2+y^2+x+y=0
\end{equation}\\
Q.12 Suppose a,b denote the distinct real roots of the quadratic polynomial $x^2+20x-2020$ and suppose c,d denote the distinct complex roots of the quadratic polynomial $x^2-20x+2020$. Then the value of
$$ac(a-c)+ad(a-d)+bc(b-c)+bd(b-d)$$ is\\
(A)0\\
(B)8000\\
(C)8080\\
(D)16000\\
                      $$SECTION2$$\\ 
\begin{center}
\begin{tikzpicture}
\node[draw]{\textbullet The section contains 2 question each question for +4 marks};
\end{tikzpicture} \\
\end{center}                 
Q.13 Let $f$: R → R  be the functions defined by
$$f(x)=\frac{x^2-3x-6}{x^2+2x+4}$$
Then which of the statement is(are) true ?\\
(A)$f$ is decreasing in the interval (-2,-1)\\
(B)$f$ is increasing in the interval (1,2)\\
(C)$f$ is onto\\
(D)Range of $f$ is $[\frac{-3}{2},2]$\\
\\
Q.14 The value of integral 
$$
\int_{0}^{\frac{\pi}{2}} \frac{3\sqrt{\cos\theta}}{(\sqrt{\cos\theta}+\sqrt{\sin \theta})} d\theta
$$
equals?\\
\newpage
                     $$SECTION3$$\\
\begin{center}
\begin{tikzpicture}
\node[draw]{\textbullet The section contains 2 question each question for +3 marks};
\end{tikzpicture} \\
\end{center}
Q.15 Let $a_1,a_2,a_2,\dots $ be a sequence of positive integers in arithmetic progression with common difference 2. Also, let $a_1,a_2,a_2,\dots $ be a sequence of positive integers in geometric progression with common ratio 2. If $a_1=b_1=c$ then the number of all possible values of c, for which the equality$$2(a_1+a_2+\dots +a_n)=b_1+b_2+\dots +b_n$$
holds for some positive integer n is ?                    
\end{document}